\title{Measuring the Impact of Early-2025 AI on Experienced Open-Source Developer Productivity}
\begin{document}
\maketitle

\begin{abstract}
Despite widespread adoption, the impact of AI tools on software development in the wild remains understudied. We conduct a randomized controlled trial (RCT) to understand how AI tools at the February–June 2025 frontier affect the produc- tivity of experienced open-source developers. 16 developers with moderate AI experience complete 246 tasks in mature projects on which they have an aver- age of 5 years of prior experience. Each task is randomly assigned to allow or disallow usage of early-2025 AI tools. When AI tools are allowed, developers primarily use Cursor Pro, a popular code editor, and Claude 3.5/3.7 Sonnet. Be- fore starting tasks, developers forecast that allowing AI will reduce completion time by 24\%. After completing the study, developers estimate that allowing AI reduced completion time by 20\%. Surprisingly, we find that allowing AI actually increases completion time by 19\%—AI tooling slowed developers down. This slowdown also contradicts predictions from experts in economics (39\% shorter) and ML (38\% shorter). To understand this result, we collect and evaluate evi- dence for 21 properties of our setting that a priori could contribute to the observed slowdown effect—for example, the size and quality standards of projects, or prior developer experience with AI tooling. Although the influence of experimental ar- tifacts cannot be entirely ruled out, the robustness of the slowdown effect across our analyses suggests it is unlikely to primarily be a function of our e
\end{abstract}

\section{Recovered Text}
Measuring the Impact of Early-2025 AI on
Experienced Open-Source Developer Productivity
Joel Becker∗, Nate Rush∗, Beth Barnes, David Rein
Model Evaluation \& Threat Research (METR)
Abstract
Despite widespread adoption, the impact of AI tools on software development in
the wild remains understudied. We conduct a randomized controlled trial (RCT)
to understand how AI tools at the February–June 2025 frontier affect the produc-
tivity of experienced open-source developers. 16 developers with moderate AI
experience complete 246 tasks in mature projects on which they have an aver-
age of 5 years of prior experience. Each task is randomly assigned to allow or
disallow usage of early-2025 AI tools. When AI tools are allowed, developers
primarily use Cursor Pro, a popular code editor, and Claude 3.5/3.7 Sonnet. Be-
fore starting tasks, developers forecast that allowing AI will reduce completion
time by 24\%. After completing the study, developers estimate that allowing AI
reduced completion time by 20\%. Surprisingly, we find that allowing AI actually
increases completion time by 19\%—AI tooling slowed developers down. This
slowdown also contradicts predictions from experts in economics (39\% shorter)
and ML (38\% shorter). To understand this result, we collect and evaluate evi-
dence for 21 properties of our setting that a priori could contribute to the observed
slowdown effect—for example, the size and quality standards of projects, or prior
developer experience with AI tooling. Although the influence of experimental ar-
tifacts cannot be entirely ruled out, the robustness of the slowdown effect across
our analyses suggests it is unlikely to primarily be a function of our experimental
design.
1
Introduction
Software development is an important part of the modern economy, and a key domain for under-
standing and forecasting AI capabilities [1; 2]. Frontier AI systems demonstrate impressive capabil-
ities on a wide range of software benchmarks [3; 4; 5; 6; 7; 8; 9] and in experiments measuring AI’s
impact on developer productivity when completing synthetic tasks [10; 11]. However, tasks used in
these lab experiments sacrifice realism for scale and efficiency: the tasks are typically self-contained,
do not require much prior context/familiarity to understand and complete, and use algorithmic eval-
uation metrics which do not capture many important capabilities [12; 13; 14]. As a result, it can be
difficult to draw inferences from results on these evaluations about AI’s impact in practice.
To reduce the inferential gap between measurements of AI capabilities and real-world impact, one
can measure the impact of AI systems in real-world settings (i.e. field experiments). Existing field
experiments aimed at measuring AI’s impact on software development measure outcomes like num-
ber of added lines of code or number of tasks completed [15; 16; 17]. However, AI systems can
affect these outcomes without productivity actually increasing—for example, code can be more ver-
bose but functionally equivalent, and tasks can be broken up into multiple smaller tasks without the
total amount of work changing—making it challenging to interpret these results.
*Equal contribution. Correspondence to \{nate, joel\}@metr.org
arXiv:2507.09089v2  [cs.AI]  25 Jul 2025
Figure 1: Experts and study participants (experienced open-source contributors) substantially over-
estimate how much AI assistance will speed up developers—tasks take 19\% more time when study
participants can use AI tools like Cursor Pro. See Appendix D for detail on speedup percentage and
confidence interval methodology.
To directly measure the impact of AI tools on developer productivity, we conduct a randomized
controlled trial by having 16 developers complete 246 tasks (2.0 hours on average) on well-known
open-source repositories (23,000 stars on average) they regularly contribute to. Each task is ran-
domly assigned to allow or disallow AI usage, and we measure how long it takes developers to
complete tasks in each condition1. Developers, who typically have tens to hundreds of hours of
prior experience using LLMs2, use AI tools considered state-of-the-art during February–June 2025
(primarily Cursor Pro with Claude 3.5/3.7 Sonnet). We collect screen recordings as they work,
providing a rich data source for analysis.
Before tasks are randomized, developers forecast that allowing AI will reduce completion time by
24\%. After study participation, developers estimate that allowing AI reduced completion time by
20\%. Surprisingly, we find that allowing AI actually increases completion time by 19\%— devel-
opers are slower when using AI tooling. Figure 1 displays this observed slowdown in contrast with
forecasts and post-hoc developer estimates of speedup from AI. We also collect forecasts of speedup
from machine learning and economics experts in academia and industry, and find that they also sub-
stantially overestimate our observed speedup.
To understand this surprising result, we manually label 143 hours of recordings of developers’ com-
puter screens while they work (representing 29\% of the total hours spent by developers), which
allows us to decompose how they spend their time when working with and without AI assistance
at a resolution of ∼10 seconds. We additionally collect rich statistics from source-code manage-
ment systems, interview and survey participating developers, and conduct subset analyses to better
understand the nature of the slowdown result.
Using these various sources of data, we identify 21 properties of our setting and experimental de-
sign that we hypothesize a priori may contribute to the slowdown effect. We group these factors
into four categories: a) direct productivity loss, b) experimental artifact, c) factors raising human
performance, and d) factors limiting AI performance. We find evidence that 5 factors contribute to
the slowdown effect, we find mixed/unclear/no evidence for 10 factors, and we find evidence against
6 factors contributing to the slowdown effect. Section 3.3 presents these factors at a high level, and
Appendix C discusses each factor in detail. While we can’t completely rule out the impact of exper-
1Crucially, the tasks are defined before they are randomized, limiting the impact of effects from AI assistance
unrelated to productivity (e.g., more verbose but functionally equivalent code)
2While 93\% of developers have previously used LLMs, only 44\% have prior experience using the Cursor
IDE.
2
imental artifacts, the slowdown effect appears broadly robust across a wide range of experimental
design decisions.
That said, many of the factors we find evidence for contributing to slowdown are specific to the
setting we study—these results do not imply that current AI systems are not useful in many realistic,
economically relevant settings. Furthermore, these results do not imply that future models will not
speed up developers in this exact setting—this is a salient possibility given the rapid pace of progress
in AI capabilities recently [2]. Finally, it remains possible that further improvements to current AI
systems (e.g. better prompting/agent scaffolding, or domain-specific finetuning) could yield positive
speedup in this setting.
Nonetheless, our results reveal a large disconnect between perceived and actual AI impact on de-
veloper productivity. Despite widespread adoption of AI tools and confident predictions of positive
speedup from both experts and developers, we observe that AI actually slows down experienced
developers in this setting.
1.1
Background
Speedup, but on synthetic tasks
Literature on productivity improvements on software tasks due
to AI usage broadly finds that AI tools increase productivity. Peng et al. [10] and Paradis et al. [11]
find 56\% and 21\% speedups on coding tasks when using AI assistance, and Weber et al. [18] finds
a 65\% increase in the rate of task requirements satisfied with AI tools. However, these studies use
artificial/synthetic tasks that make it difficult to directly draw inferences about the real-world impact
of AI tools. For example, Peng et al. [10] asks developers to implement a very basic HTTP server
in JavaScript to satisfy several automatic test cases that are shown to the developers—this task is a)
unrepresentative of most software development work, and b) likely to be similar to a large amount
of LLM training data, which may unfairly advantage AI systems relative to humans.
Speedup, but with non-fixed outcome measures
Other literature uses tasks found “in the wild,”
either via natural experiments [16] or randomized controlled trials [15; 17], finding 14-51\% in-
creases in output productivity metrics. However, these studies use outcome measures that are not
fixed in advance—i.e. lines of code written, number of code commits, and pull requests3 (PRs)
as their key outcome measures respectively. It’s possible for AI assistance to affect the outcomes
without actually increasing productivity, e.g. by causing developers to write more verbose but func-
tionally equivalent code, or causing them to break up pull requests into smaller chunks of work.
Impressive AI benchmark results
This general consensus around AI tooling’s effect on software
developer productivity is perhaps unsurprising, given the impressive apparent capabilities of frontier
AIs on challenging question-answering and agentic tasks used in popular AI benchmarks [19; 20].
Heterogeneous effects by experience
One important question that emerges given these impres-
sive results is whether productivity gains are captured by individuals of all experience levels. The
canonical framework of Agrawal et al. [21] treats AI as a fall in the cost of prediction, with distri-
butional consequences depending on which complementary sub-problems the tool does not solve.
Existing empirical work on the micro-level effects of generative AI tools tends to find that ac-
cess to these tools benefits less experienced workers more, compressing performance distributions
[22; 23; 10; 24].
These heterogeneous effects motivate our focus on highly skilled open-source developers, as there
has been relatively less research in this setting.
Mixed speedup results in other domains
Some literature measures the impact of frontier AI
systems in settings other than software development, for example, for CBRN uplift risk assessment,
finding mixed results with recent AI systems [25; 26; 27; 28]. Other research finds substantial
productivity increases in non-software domains [22; 23].
Understanding AI’s economic impact
Finally, some literature tries to predict how AI capability
advances might a) affect the rate of AI progress (e.g. if AI systems can substitute for human AI
3See Appendix F for a primer on open-source development terminology.
3
Figure 2: Our experimental design. Tasks (referred to as issues) are defined before treatment as-
signment, screen recordings let us verify compliance (and provide a rich source of data for analysis),
and forecasts from experts and developers help us measure the gap between expectations and ob-
served results.
R\&D labor), or b) broadly impact the economy. Leibowich et al. [29] interview AI researchers
about how full automation of AI R\&D might alter the pace of advancement, several papers explore
the possibility of explosive economic growth via large-scale AI labor substitution [30; 31; 32], and
the economics literature includes both optimistic and skeptical perspectives on AI’s productivity
impact [33; 34; 35].
Our study primarily complements existing literature measuring the impact of AI on software devel-
opment by:
1. Testing AI models at the February–June 2025 frontier,
2. Using unfiltered, “live” open-source repository tasks rather than synthetic or cherry-picked
tasks,
3. Using a fixed outcome measure (speedup on tasks defined before randomized treatment
assignment),
4. Recruiting experienced engineers with years of expertise in the target repositories, and
5. Collecting rich data on time usage, AI code suggestions, and developers’ qualitative expe-
riences.
2
Methodology
2.1
Developers and Repositories
We recruit experienced developers from large open source repositories to work on real tasks defined
on these repositories. Developers come from a mix of our professional networks and from outreach
to active contributors to large, popular Github repositories. The developers are experienced software
engineers (typically over a decade of experience), and are regular contributors to the repositories
we use—on average, they have 5 years of experience working on their repository, representing 59\%
of that repository’s lifetime, over which time they have made 1,500 commits to the repo. As an
incentive to participate, we pay developers \$150/hour. Appendix G provides more detail about our
recruitment and incentivization process.
The repositories themselves are large and mature. On average, they have 23,000 stars, 1,100,000
lines of code, 4,900 forks, 20,000 commits, and 710 committers, and they broadly have very high
quality bars for code contributions. For example, one set of repository contribution guidelines con-
cludes: “Phew. While the above may be a lot to remember [..] the motivation for enforcing process
is to ensure that all code contributions meet a certain quality threshold.” Section G.7 details further
statistics about individual developers and repositories.
4
Figure 3: Real issues completed during the study from the stdlibjs and mito repositories
2.2
Experimental Design
Each developer provides a list of real issues in their repository to work on as part of this study.
Issues are typically bug reports, feature requests, or work items used to coordinate development.
They range from brief problem descriptions to detailed analyses and represent work ranging from
minutes to hours. Two example issues are shown in Figure 3. Many issues are defined before the
study period begins, but some are created during the study period.4
After collecting this issue list, developers forecast how long each issue would take if they were
to complete it both with and without AI assistance. We use these forecasts as a proxy for issue
difficulty, and to measure per-issue speedup anticipated by the developer. These issues are then
randomized to one or the other condition via a simulated fair coin flip.5 If AI is allowed, developers
can use any AI tools or models they choose, including no AI tooling if they expect it to not be helpful.
If AI is not allowed, no generative AI tooling can be used.6
Developers then work on their assigned issues in their preferred order—they are allowed to flexibly
complete their work as they normally would, and sometimes work on multiple issues at a time. After
completing an issue to their satisfaction, they submit a pull request (PR) to their repository, which
is typically reviewed by another developer. They make any changes suggested by the PR reviewer,
4About half of issues included in the study were not formally defined on the repository, and instead were
tracked separately for our experiment. Importantly, all issues represent real work that developers wanted to
contribute to their repositories, and all work completed by developers is submitted and reviewed through each
repository’s standard source-code management system (e.g. GitHub/GitLab). Developers are asked to con-
tribute issues taking a maximum of two hours, or to break up issues taking longer into multiple issues.
525 issues early in the study were randomized differently. We show that results are not sensitive to inclu-
sion/exclusion of these issues in Section C.3.4, and we include these issues for statistical power.
6AI-based tab autocomplete is disallowed in the AI-disallowed condition if it uses LLMs (e.g. GitHub
Copilot) but allowed otherwise. Search engines, which sometimes use AI under the hood, remain allowed in
the AI-disallowed condition.
5
and merge their completed PR into the repository7. As the repositories included in the study have
very high quality and review standards, merged PRs rarely contain mistakes or flaws. Finally, they
self-report how long they spend working on each issue before and after PR review.
See Section G.2 for the full written instructions given to developers before they start working.
2.2.1
AI Tools and Training
Two popular means of using modern large language model (LLM) based AI tools are via web-based
user interfaces (e.g. chatgpt.com) and the integrated development environment (IDE) Cursor Pro
(which we provide a subscription for). Cursor is a fork of the widely used VSCode IDE with near-
identical features, that additionally includes extra AI features like a language model chat interface,
and an AI agent tool that can search and edit files, run arbitrary bash commands, prompt/ask the
user for more details when relevant, and iterate/debug programs without constant input from users.
Developers have a range of experience using AI tools: 93\% have prior experience with tools like
ChatGPT, but only 44\% have experience using Cursor.
We provide developers with Cursor Pro subscriptions and conduct live basic training, validating that
developers are able to prompt Cursor effectively to edit files in their own codebase, accept changes,
and revert to previous checkpoints. However, we don’t require that they use Cursor specifically.
Developers working on issues for which AI is allowed can use any AI tools of their choosing, or no
AI tools if they prefer. See Section F.2 for further information on these two methods of accessing
AI assistance, and Appendix G for more detail about our training and onboarding process.
2.2.2
Data Collection
Contributors completed issues largely as they would outside of our experiment, with a few excep-
tions: they typically record their screen as they work (providing us with a source of rich data on their
AI usage and working patterns), when using AI they often use the Cursor IDE, which sometimes dif-
fers from their normal development environment (e.g. neovim, PyCharm), and they record details of
their working patterns after each issue is completed. Broadly, developers reported that they were not
significantly inconvenienced or affected by these differences compared to their normal workflows
(see section Section 3.3 for discussion of how experimental design may affect results). Developers
self-report the length of time they spend working on each issue before and after PR review.
2.3
Effect Estimation
We aim to measure the percentage change in issue time to completion due to AI, defined as:
S =
E[completion time with AI allowed]
E[completion time with AI disallowed] −1
A value of -50\% would indicate issues assigned to AI-allowed taking half the time of issues assigned
to AI-disallowed; 0\% would indicate that issues take the same time to complete with and without
AI-allowed; 100\% would indicate that issues assigned to AI-allowed take twice as long as issues
assigned to AI-disallowed. For convenience and clarity, we typically report the absolute value of S,
and colloquially refer to negative values of as speedup, and positive values as slowdown8
We use the total implementation time (including both pre- and post-pull request review time) devel-
opers self-report spending on each issue to estimate the percentage change in issue completion time
due to AI. A small number of issues do not end up getting reviewed/merged for various reasons—
we impute the post-review implementation time of these issues with the mean post-review time
per treatment condition. Section C.3.4 shows that our results are robust to the specific imputation
method used.
We estimate S via a log-linear regression for total implementation time, incorporating estimates of
issue difficulty to reduce noise in our estimate. We also use adapted versions of this regression when
7Not all PRs end up getting merged by the end of the study period—see Section C.3.4 for discussion of how
this may affect results.
8This is a mild abuse of language, because speedup and slowdown refer to rates, while the quantity we
actually estimate and report is the percentage change in time from AI.
6
examining how performance varies in different subsets of the data (to efficiently account for any
post-randomization issue difficulty imbalances that may be more pronounced within subsets). See
Appendix D for details of the regression approach used.
2.4
Supplementary Data Collection
Developer Forecasts and Post hoc Estimates
Before randomizing each issue to either AI-allowed
or AI-disallowed, we ask developers to forecast how long they expect each issue to take if they are
allowed to use AI assistance, and if they are disallowed from using AI. This provides a measure
of issue difficulty independent of random treatment assignment, and enables us to assess developer
calibration and expectations regarding the usefulness of AI assistance for their work.
After the study’s conclusion, we ask developers to estimate post hoc how much AI tooling affected
the time it took them to complete issues. This helps us understand the relationship between individ-
uals’ beliefs about AI helpfulness and our measured effects.
The detailed instructions given to developers for per-issue forecasts and exit surveys are available in
Section G.2.
Expert Forecasts
We solicit forecasts from experts in machine learning and economics in order
to compare our results against expert prior expectations regarding AI capabilities [36]. Machine
learning experts are a mix of industry professionals from frontier AI labs, and academics from
universities with leading graduate computer science programs. Economics experts are primarily
academics and PhD students from universities with leading economics graduate programs. See
Section E.5 for summary statistics of expert forecasts by expert group.
Experts are given detailed information about our study design, repositories participating in our ex-
periment, and the proportion of developers who have prior experience with AI tooling (Section G.9).
They then report their prediction of our point estimate for E[time with AI disallowed]
E[time with AI allowed]
9.
To incentivize experts to make accurate forecasts, we offer to pay max(0, \$50×(1−(true answer−
guessed answer)2)) for their point estimates.1011
Cursor Analytics Data
We collect per-developer analytics data from Cursor’s analytics dash-
board, for the 13 developers we provide Cursor Pro subscriptions to. For each developer, we collect
their most-used model and the number of suggested and accepted lines of AI-generated code.
Exit Interviews and Surveys
After the study’s conclusion, we conduct exit interviews and surveys
with all developers to assess where they found AI helpful, what strategies they used to effectively
elicit work from AI, whether they felt they improved at using AI tooling over the course of the
study, and to estimate how much they were sped up by AI during the study. Full details of the exit
interviews are available in Section G.5.1.
Qualitative Evidence
Throughout the study we collect qualitative evidence from developers, to
form a more comprehensive understanding of their experiences using AI tools. Developers are in-
structed to take detailed notes regarding their experience and usage of AI tools, and we use inductive
coding—where we iteratively read through the data to identify recurring patterns, create and refine
categories as they emerge, and reorganize excerpts until stable themes develop—to cluster excerpts
from these notes. Qualitative results in Section 3.3 (particularly quotes from developers) were col-
lected primarily with this methodology. However, to investigate initial hypotheses, we often ask
developers probing/targeted questions, so we cannot rule out bias from this type of qualitative evi-
dence.
9Note that the estimate we use for Figure 1 is transformed to represent E[T | AI=1]
E[T | AI=0] −1. It is not necessarily
true that forecasters’ belief about the point estimate of E[T | AI=1]
E[T | AI=0] is equal to the reciprocal of their belief about
the point estimate of E[T | AI=0]
E[T | AI=1].
10Approximately one-third of forecasters are offered a maximum of \$100 rather than \$50.
11Taking a maximum with 0 makes our scoring rule improper [37]—their reward is not necessarily maxi-
mized at their true belief. We use this scoring rule for simplicity and clarity.
7
Figure 4: Left: Raw average forecasted implementation times. Right: Raw average observed im-
plementation times. The ratio of observed implementation times gives a more extreme slowdown
estimate than regression-based estimates (Section D.1) because AI-allowed issues are forecasted
(importantly, before treatment assignment) by developers to take slightly longer, which the regres-
sion corrects for. Both: Section D.5 describes confidence intervals around ratios of average times.
Fine-Grained Screen Recording Activity Labels
To compare how developers spend their time
with and without AI assistance, we manually label a subset of 128 screen recordings with fine-
grained activity labels, totaling 143 hours of video. In results based on these labeled screen record-
ings, we filter to remove issues where we find cheating, issues where the screen recording is broken
for >10\% of recording time, and issues with a >20\% discrepancy between self-reported time and
the recording time. This results in 74 recordings representing 84 hours of video for further analysis.
We label whether developers are: actively writing code, testing and debugging their code, reading
or searching for information, using Git or managing their environment, prompting an AI system,
waiting on an AI system to generate output, reviewing AI outputs, or idling/doing other miscella-
neous work. Each high-level label is further broken down into one of 27 fine-grained categories.
Labels have a resolution of ∼10 seconds. Section G.8 describes the instructions and process used
for labeling screen recordings.
3
Results
Developers complete 136 issues with AI-allowed and 110 issues with AI-disallowed. Section G.7
shows the number of issues completed across repositories and developers, respectively. We find that
when developers use AI tools, they implement issues in 19\% more time on average (Figure 1), and
nearly all quantiles of observed implementation time see AI-allowed issues taking longer (Figure 5).
That is, developers are slower when using AI is allowed. Colloquially, we refer to this result that
issues with AI-allowed take longer than issues with AI-disallowed as slowdown.
3.1
Forecasts
Developer Forecasts and Post hoc Estimates
Before developers complete each issue, they fore-
cast how long they expect them to take with and without AI assistance. On average, they forecast
speedup of 24\%. Interestingly, after the experiment they post-hoc estimate that they were sped-up
by 20\% when using AI is allowed—after they used AI assistance, they estimate similar speedup as
8
Figure 5: Empirical cumulative distribution functions of observed implementation times. Percentile
ordering is calculated separately for each treatment group.
before, despite the fact that they are in fact slowed down by 19\% (Figure 1). Figure 4 displays the
raw average forecasted and observed implementation times12.
Despite developers forecasting speedup from AI (while they are slowed down), developer forecasts
are informative about completion time—the Pearson correlation between the time developers fore-
cast AI-allowed issues taking and the actual time they take is 0.64, and the correlation between the
time developers forecast AI-disallowed issues taking and the actual time they take is 0.59. This
suggests that developers are broadly well-calibrated on the relative amount of time that issues will
take, but their expectations regarding the usefulness of AI assistance are reversed.
Expert Forecasts
Speedup forecasts from 34 economics experts and 54 machine learning experts
overestimate speedup even more drastically than developers, predicting AI will lead to decreases
in implementation time of 39\% and 38\%, respectively (Figure 1). We show distributions of expert
forecasts in Section E.5.
3.2
Activity Labels
On a subset of 74 issues for which we have valid screen recordings, we manually label the activities
developers engage in while they work. Figure 6 shows the percentage of time developers spend
for each type of issue (AI-allowed or AI-disallowed). When allowed to use AI, developers spend a
smaller proportion of their time actively coding and reading/searching for information. Instead, they
spend time reviewing AI outputs, prompting AI systems, and waiting for AI generations. Interest-
ingly, they also spend a somewhat higher proportion of their time idle, where their screen recording
doesn’t show any activity. Section E.4 shows the number of minutes spent on average in each cate-
gory (instead of percentage time spent), as well as the distributions of labels broken down into more
fine-grained activities.
12The raw percentage difference in implementation times between AI-allowed and AI-disallowed issues is
34\%, which is larger in absolute value than the 19\% slowdown estimated using the regression specified in
Section D.1. This is true because AI-allowed issues ended up being slightly more difficult than AI-disallowed
issues after randomization, as measured by the forecasted AI-disallowed times (not statistically significant; see
Table 4). Our regression accounts for this, while this simple ratio estimator does not. See Figure 13 for results
implied by alternative estimators.
9
Figure 6: On the subset of labeled screen recordings, when AI is allowed, developers spend less
time actively coding and searching for/reading information, and instead spend time prompting AI,
waiting on and reviewing AI outputs, and idle. Figure 19 shows the absolute (average) minutes spent
in each category, and Figure 21 presents these results broken down into 27 fine-grained categories.
3.3
Factor Analysis
Given the surprising nature of this result, we investigate 21 potential contributing factors that may
contribute to developers spending more time on tasks when AI usage is allowed. We group these
factors into four categories:
• Direct productivity loss (Ý): mechanisms by which the use of AI tools actively slows
down development.
• Experimental artifact (e): confounders from our experimental setup or procedures that
may introduce biases, or limit the external validity.
• Raises developer performance ( ): attributes of the issues, repositories, or setting that
improve developer ability relative to AI.
• Limits AI performance (Æ): attributes of the issues, repositories, or AI/environment tool-
ing that diminish AI’s effectiveness relative to developers.
Using entry and exit surveys, screen recordings, developer interviews, and subset analyses we
find qualitative and quantitative evidence that 5 of the 21 factors contribute to slowdown, we find
mixed/unclear/no evidence that 10 of the factors contribute to slowdown, and we find evidence
against 6 of the factors contributing. However, we strongly caution against over-indexing on the ba-
sis of any individual pieces of evidence, as we are not powered for statistically significant multiple
comparisons when subsetting our data. This analysis is intended to provide speculative, suggestive
evidence about the mechanisms behind slowdown. Appendix C discusses the evidence for/against
each factor in Table 1.
4
Discussion
We provide evidence that recent AI systems slow down experienced open-source developers with
moderate AI experience completing real issues on large, popular repositories they are highly familiar
with. This observed slowdown serves as some evidence that AI capabilities in the wild may be lower
than results on commonly used benchmarks may suggest.
Furthermore, we show that both experts and developers drastically overestimate the usefulness of AI
on developer productivity, even after they have spent many hours using the tools. This underscores
10
Factors likely to contribute to slowdown
Factor
Type
Relevant Observations
Over-optimism
about
AI
usefulness
(C.1.1)
Ý
• Developers forecast AI will decrease implementation time by 24\%
• Developers post hoc estimate AI decreased implementation time by 20\%
High developer familiarity with reposito-
ries (C.1.2)
• Developers slowed down more on issues they are more familiar with
• Developers report that their experience makes it difficult for AI to help them
• Developers average 5 years experience and 1,500 commits on repositories
Large and complex repositories (C.1.3)
Æ
• Developers report AI performs worse in large and complex environments
• Repositories average 10 years old with >1,100,000 lines of code
Low AI reliability (C.1.4)
Æ
• Developers accept <44\% of AI generations
• Majority report making major changes to clean up AI code
• 9\% of time spent reviewing/cleaning AI outputs
Implicit repository context (C.1.5)
Æ
• Developers report AI doesn’t utilize important tacit knowledge or context
Factors with unclear effect on slowdown
Factor
Type
Relevant Observations
Experimentally driven overuse of AI
(C.2.1)
e
• Developers sometimes report overuse due to experiment
• Similar slowdown from developers reporting overuse vs. normal use
Unrepresentative
task
distribution
(C.2.2)
e
• Developers report issues are standard but on the shorter side
• Excludes non-programming tasks developers complete in normal work
AI increasing issue scope (C.2.3)
e
• Developers who report scope creep see less slowdown
• Mixed developer reports on AI’s impact on scope
• 47\% more lines of code per forecasted hour in AI-allowed issues
Bias from issue completion order (C.2.4)
e
• Developers decide order post randomization
Sampling bias in developer recruitment
(C.2.5)
e
• Developers who rely heavily on AI may be less likely to participate
Trading speed for ease (C.2.6)
Ý
• Some developers report using AI is less effortful
• High developer retention on Cursor
Low quality initial pull requests (C.2.7)
Ý
• Minor difference in mean post-review times between conditions
• Qualitatively similar PR quality between conditions
Below-average use of AI tools (C.2.8)
Æ
• Similar slowdown for developers with prior Cursor experience
• No clear learning effect across first 30-50 hours of Cursor usage
• Developers appear qualitatively in distribution for Cursor Pro users
AI generation latency (C.2.9)
Æ
• Mixed developer reports that waiting on AI generations was important
• Developers spend 4\% of time waiting on AI generations
Suboptimal elicitation (C.2.10)
Æ
• Developers use Cursor agents/chat in majority of AI-allowed issues
• Developers sample few tokens from models
• But existing literature finding positive speedup also uses few tokens
• Unused elicitation strategies could improve AI reliability
Factors unlikely to contribute to slowdown
Factor
Type
Relevant Observations
Unfamiliar
development
environment
(C.3.1)
e
• Most developers use comparable IDEs between treatment conditions
• These developers still see slowdown of 24\%
• No clear learning effects across first 30-50 hours of Cursor usage
Cheating or under-use of AI (C.3.2)
e
• AI used in all but 16.4\% of allowed cases with labeled screen recordings
• Only 3 cheating instances in 54 screen recordings
Issue dropout (C.3.3)
e
• Developers with no accidental dropout see similar slowdown
• Issues dropped intentionally are qualitatively unbiased
Non-robust outcome measure (C.3.4)
e
• Alternative outcome measures yield similar slowdown
Non-robust estimator (C.3.5)
e
• Alternative estimators yield similar slowdown
Non-frontier model usage (C.3.6)
Æ
• Developers mostly use (at the time) frontier models
Table 1: Summary of factors that may a priori explain or contribute to slowdown, grouped by the
state of evidence for or against their impact on the slowdown effect.   are factors that raise human
performance, Æ are factors that limit AI performance, e are experimental artifacts that may bias/-
confound results, and Ý are factors that directly contribute to productivity losses.
11
the importance of conducting field experiments with robust outcome measures, compared to relying
solely on expert forecasts or developer surveys.
4.1
Key Caveats
Setting-specific factors
We caution readers against overgeneralizing on the basis of our results.
The slowdown we observe does not imply that current AI tools do not often improve developer’s
productivity—we find evidence that the high developer familiarity with repositories and the size and
maturity of the repositories both contribute to the observed slowdown, and these factors do not apply
in many software development settings. For example, our results are consistent with small greenfield
projects or development in unfamiliar codebases seeing substantial speedup from AI assistance.
AI-specific factors
We expect that AI systems that have higher fundamental reliability, lower
latency, and/or are better elicited (e.g. via more inference compute/tokens, more skilled prompt-
ing/scaffolding, or explicit fine-tuning on repositories) could speed up developers in our setting (i.e.
experienced open-source developers on large repositories).
Agents can make meaningful progress on issues
We have preliminary evidence (forthcoming)
that fully autonomous AI agents using Claude 3.7 Sonnet can often correctly implement the core
functionality of issues on several repositories that are included in our study, although they fail to
fully satisfy all requirements (typically leaving out important documentation, failing linting/styling
rules, and leaving out key unit or integration tests). This represents immense progress relative to the
state of AI just 1-2 years ago, and if progress continues apace (which is a priori at least plausible,
although not guaranteed), we may soon see significant speedup in this setting.
5
Acknowledgments
We thank the open-source developers who participated in this study. Your hard work, diligent
record keeping, and excellent software made it a pleasure to work with you. Thanks to Aaron
Diamond-Reivich, Alan Akbik, Domenic Denicola, Dens Sumesh, Jaden Fiotto-Kaufman, Jo˜ao
Gante, Liam DeVoe, Matthew Pickering, Muhammad Haris, Philipp Burckhardt, Quentin Anthony,
Ruben Bloom, Sam Derbyshire, and other participating developers.
We thank the following reviewers for feedback on the experimental design and paper drafts:
Adrien Ecoffet, Alexander Barry, Ali Merali, Ajeya Cotra, Andres Campero, Andrey Fradkin, Basil
Halperin, Cozmin Ududec, Eli Lifland, Ernest Davis, Gregory Sun, Hjalmar Wijk, James Requeima,
Jide Alaga, Josh Jacobson, Lawrence Chan, Megan Kinniment, Michael Sklar, Neev Parikh, Rif A.
Saurous, Rob Miles, Ryan Greenblatt, Seraphina Nix, Sydney Von Arx, Thomas Kwa, and Tom
Cunningham.
We thank the following individuals for help with data collection: Adam Hanson, Amy Ngo, Chris
Canal, Jebastin Nadar, Luis Slyfield, and Martin Milbradt.
We thank Sami Jawar and Thomas Broadley for technical support throughout the project.
We thank the following for their operational support through the project: Bhaskar Chaturvedi, Emma
Abele, Kit Harris, Kris Chari, Kyle Scott, Rebecca Baron, and Rae She.
The authors thank Stephanie He for graphic design contributions.
The authors especially thank Aron Lajko, Chris Painter, Jasmine Dhaliwal, and Steve Newman for
close review, feedback, and support throughout the project.
12
References
[1] Nestor Maslej, Loredana Fattorini, Raymond Perrault, Yolanda Gil, Vanessa Parli, Njenga Kar-
iuki, Emily Capstick, Anka Reuel, Erik Brynjolfsson, John Etchemendy, Katrina Ligett, Terah
Lyons, James Manyika, Juan Carlos Niebles, Yoav Shoham, Russell Wald, Tobi Walsh, Armin
Hamrah, Lapo Santarlasci, Julia Betts Lotufo, Alexandra Rome, Andrew Shi, and Sukrut
Oak. Artificial intelligence index report 2025, 2025. URL https://arxiv.org/abs/2504.
07139.
[2] Thomas Kwa, Ben West, Joel Becker, Amy Deng, Katharyn Garcia, Max Hasin, Sami Jawhar,
Megan Kinniment, Nate Rush, Sydney Von Arx, Ryan Bloom, Thomas Broadley, Haoxing Du,
Brian Goodrich, Nikola Jurkovic, Luke Harold Miles, Seraphina Nix, Tao Lin, Neev Parikh,
David Rein, Lucas Jun Koba Sato, Hjalmar Wijk, Daniel M. Ziegler, Elizabeth Barnes, and
Lawrence Chan. Measuring ai ability to complete long tasks, 2025. URL https://arxiv.
org/abs/2503.14499.
[3] Hjalmar Wijk, Tao Lin, Joel Becker, Sami Jawhar, Neev Parikh, Thomas Broadley, Lawrence
Chan, Michael Chen, Josh Clymer, Jai Dhyani, Elena Ericheva, Katharyn Garcia, Brian
Goodrich, Nikola Jurkovic, Holden Karnofsky, Megan Kinniment, Aron Lajko, Seraphina Nix,
Lucas Sato, William Saunders, Maksym Taran, Ben West, and Elizabeth Barnes. Re-bench:
Evaluating frontier ai r\&d capabilities of language model agents against human experts, 2025.
URL https://arxiv.org/abs/2411.15114.
[4] Jun Shern Chan, Neil Chowdhury, Oliver Jaffe, James Aung, Dane Sherburn, Evan Mays,
Giulio Starace, Kevin Liu, Leon Maksin, Tejal Patwardhan, Lilian Weng, and Aleksander
Madry. Mle-bench: Evaluating machine learning agents on machine learning engineering,
2025. URL https://arxiv.org/abs/2410.07095.
[5] Giulio Starace, Oliver Jaffe, Dane Sherburn, James Aung, Jun Shern Chan, Leon Maksin,
Rachel Dias, Evan Mays, Benjamin Kinsella, Wyatt Thompson, Johannes Heidecke, Amelia
Glaese, and Tejal Patwardhan. Paperbench: Evaluating ai’s ability to replicate ai research,
2025. URL https://arxiv.org/abs/2504.01848.
[6] Shanghaoran Quan, Jiaxi Yang, Bowen Yu, Bo Zheng, Dayiheng Liu, An Yang, Xuancheng
Ren, Bofei Gao, Yibo Miao, Yunlong Feng, Zekun Wang, Jian Yang, Zeyu Cui, Yang Fan,
Yichang Zhang, Binyuan Hui, and Junyang Lin. Codeelo: Benchmarking competition-level
code generation of llms with human-comparable elo ratings, 2025. URL https://arxiv.
org/abs/2501.01257.
[7] Samuel Miserendino, Michele Wang, Tejal Patwardhan, and Johannes Heidecke. Swe-lancer:
Can frontier llms earn \$1 million from real-world freelance software engineering?, 2025. URL
https://arxiv.org/abs/2502.12115.
[8] David Rein, Betty Li Hou, Asa Cooper Stickland, Jackson Petty, Richard Yuanzhe Pang, Julien
Dirani, Julian Michael, and Samuel R. Bowman. Gpqa: A graduate-level google-proof q\&a
benchmark, 2023. URL https://arxiv.org/abs/2311.12022.
[9] David Rein, Joel Becker, Amy Deng, Seraphina Nix, Chris Canal, Daniel O’Connel, Pip
Arnott, Ryan Bloom, Thomas Broadley, Katharyn Garcia, Brian Goodrich, Max Hasin, Sami
Jawhar, Megan Kinniment, Thomas Kwa, Aron Lajko, Nate Rush, Lucas Jun Koba Sato, Syd-
ney Von Arx, Ben West, Lawrence Chan, and Elizabeth Barnes. Hcast: Human-calibrated
autonomy software tasks, 2025. URL https://arxiv.org/abs/2503.17354.
[10] Sida Peng, Eirini Kalliamvakou, Peter Cihon, and Mert Demirer. The impact of ai on developer
productivity: Evidence from github copilot, 2023. URL https://arxiv.org/abs/2302.
06590.
[11] Elise Paradis, Kate Grey, Quinn Madison, Daye Nam, Andrew Macvean, Vahid Meimand,
Nan Zhang, Ben Ferrari-Church, and Satish Chandra. How much does ai impact development
speed? an enterprise-based randomized controlled trial, 2024. URL https://arxiv.org/
abs/2410.12944.
13
[12] Inioluwa Deborah Raji, Emily M. Bender, Amandalynne Paullada, Emily Denton, and Alex
Hanna. AI and the everything in the whole wide world benchmark. CoRR, abs/2111.15366,
2021. URL https://arxiv.org/abs/2111.15366.
[13] Stella Biderman, Hailey Schoelkopf, Lintang Sutawika, Leo Gao, Jonathan Tow, Baber Ab-
basi, Alham Fikri Aji, Pawan Sasanka Ammanamanchi, Sidney Black, Jordan Clive, Anthony
DiPofi, Julen Etxaniz, Benjamin Fattori, Jessica Zosa Forde, Charles Foster, Jeffrey Hsu,
Mimansa Jaiswal, Wilson Y. Lee, Haonan Li, Charles Lovering, Niklas Muennighoff, Ellie
Pavlick, Jason Phang, Aviya Skowron, Samson Tan, Xiangru Tang, Kevin A. Wang, Genta In-
dra Winata, Franc¸ois Yvon, and Andy Zou. Lessons from the trenches on reproducible evalu-
ation of language models, 2024. URL https://arxiv.org/abs/2405.14782.
[14] Ernest Davis. Benchmarks for automated commonsense reasoning: A survey, 2023. URL
https://arxiv.org/abs/2302.04752.
[15] Leonardo Gambacorta, Han Qiu, Shuo Shan, and Daniel M Rees. Generative AI and labour
productivity: a field experiment on coding, volume 1208. Bank for International Settlements,
Monetary and Economic Department, 2024.
[16] Doron Yeverechyahu, Raveesh Mayya, and Gal Oestreicher-Singer.
The impact of large
language models on open-source innovation: Evidence from github copilot, 2025.
URL
https://ssrn.com/abstract=4684662.
[17] Zheyuan Cui, Mert Demirer, Sonia Jaffe, Leon Musolff, Sida Peng, and Tobias Salz. The
effects of generative ai on high-skilled work: Evidence from three field experiments with soft-
ware developers, June 2025. URL https://ssrn.com/abstract=4945566.
[18] Thomas Weber, Maximilian Brandmaier, Albrecht Schmidt, and Sven Mayer. Significant pro-
ductivity gains through programming with large language models. Proc. ACM Hum.-Comput.
Interact., 8(EICS), June 2024. doi: 10.1145/3661145. URL https://doi.org/10.1145/
3661145.
[19] OpenAI.
OpenAI o3 and o4-mini System Card.
https://cdn.openai.com/pdf/
2221c875-02dc-4789-800b-e7758f3722c1/o3-and-o4-mini-system-card.pdf,
2025. [Accessed 23-06-2025].
[20] Anthropic.
Anthropic Claude 4 System Card.
https://www-cdn.anthropic.com/
6be99a52cb68eb70eb9572b4cafad13df32ed995.pdf, 2025. [Accessed 23-06-2025].
[21] Ajay Agrawal, Joshua S. Gans, and Avi Goldfarb. Artificial intelligence: The ambiguous labor
market impact of automating prediction. Journal of Economic Perspectives, 33(2):31–50, May
2019. doi: 10.1257/jep.33.2.31. URL https://www.aeaweb.org/articles?id=10.1257/
jep.33.2.31.
[22] Erik Brynjolfsson, Danielle Li, and Lindsey Raymond. Generative ai at work*. The Quarterly
Journal of Economics, 140(2):889–942, 02 2025. ISSN 0033-5533. doi: 10.1093/qje/qjae044.
URL https://doi.org/10.1093/qje/qjae044.
[23] Shakked Noy and Whitney Zhang. Experimental evidence on the productivity effects of gener-
ative artificial intelligence. Science, 381(6654):187–192, 2023. doi: 10.1126/science.adh2586.
URL https://www.science.org/doi/abs/10.1126/science.adh2586.
[24] Jonathan H. Choi and Daniel Schwarcz. Ai assistance in legal analysis: An empirical study.
Journal of Legal Education, 73(2), 2025.
URL https://jle.aals.org/home/vol73/
iss2/5/.
[25] Anthropic.
Responsible scaling policy evaluations report – claude 3 opus.
Technical
report, Anthropic, 2024.
URL https://cdn.sanity.io/files/4zrzovbb/website/
210523b8e11b09c704c5e185fd362fe9e648d457.pdf. Accessed June 2025.
[26] C. Mouton, Caleb Lucas, and Ella Guest.
The operational risks of ai in large-
scale
biological
attacks.
Research
report,
RAND
Corporation,
Santa
Monica,
2024.
URL https://www.rand.org/content/dam/rand/pubs/research\_reports/
RRA2900/RRA2977-2/RAND\_RRA2977-2.pdf.
14
[27] Aaron Grattafiori et al. The llama 3 herd of models, 2024. URL https://arxiv.org/abs/
2407.21783.
[28] Tejal Patwardhan, Kevin Liu, Todor Markov, Neil Chowdhury, Dillon Leet, Natalie
Cone, Caitlin Maltbie, Joost Huizinga, Carroll Wainwright, Shawn (Froggi) Jackson,
Steven Adler, Rocco Casagrande, and Aleksander Madry.
Building an early warn-
ing system for LLM-aided biological threat creation.
https://openai.com/index/
building-an-early-warning-system-for-llm-aided-biological-threat-creation/,
January 2024. Accessed: 2025-06-25.
[29] Jared Leibowich, Nikola Jurkovic, and Tom Davidson. Could advanced ai accelerate the pace
of ai progress? interviews with ai researchers, 2024. URL https://ssrn.com/abstract=
5115692.
[30] Ege Erdil and Tamay Besiroglu. Explosive growth from ai automation: A review of the argu-
ments, 2024. URL https://arxiv.org/abs/2309.11690.
[31] Ege Erdil, Andrei Potlogea, Tamay Besiroglu, Edu Roldan, Anson Ho, Jaime Sevilla, Matthew
Barnett, Matej Vrzla, and Robert Sandler. Gate: An integrated assessment model for ai au-
tomation, 2025. URL https://arxiv.org/abs/2503.04941.
[32] Tom
Davidson.
What
a
compute-centric
framework
says
about
takeoff
speeds.
https://www.openphilanthropy.org/research/
what-a-compute-centric-framework-says-about-takeoff-speeds/,
June
2023.
Open Philanthropy. Accessed: 2025-06-25.
[33] Daron Acemoglu. The simple macroeconomics of ai. Economic Policy, 40(121):13–58, 08
2024. ISSN 0266-4658. doi: 10.1093/epolic/eiae042. URL https://doi.org/10.1093/
epolic/eiae042.
[34] Ajay Agrawal, Joshua Gans, and Avi Goldfarb. Economic policy for artificial intelligence.
Innovation Policy and the Economy, 19:139–159, 2019. doi: 10.1086/699935. URL https:
//doi.org/10.1086/699935.
[35] Jason Furman and Robert Seamans.
AI and the Economy, pages 161–191.
University of
Chicago Press, May 2018. doi: 10.1086/699936. URL http://www.nber.org/chapters/
c14099.
[36] Stefano DellaVigna, Nicholas Otis, and Eva Vivalt. Forecasting the results of experiments:
Piloting an elicitation strategy. AEA Papers and Proceedings, 110:75–79, May 2020. doi: 10.
1257/pandp.20201080. URL https://www.aeaweb.org/articles?id=10.1257/pandp.
20201080.
[37] Tilmann Gneiting and Adrian E Raftery. Strictly proper scoring rules, prediction, and es-
timation.
Journal of the American Statistical Association, 102(477):359–378, 2007.
doi:
10.1198/016214506000001437. URL https://doi.org/10.1198/016214506000001437.
[38] Carlos E Jimenez, John Yang, Alexander Wettig, Shunyu Yao, Kexin Pei, Ofir Press, and
Karthik R Narasimhan. SWE-bench: Can language models resolve real-world github issues?
In The Twelfth International Conference on Learning Representations, 2024. URL https:
//openreview.net/forum?id=VTF8yNQM66.
[39] Nelson F. Liu, Kevin Lin, John Hewitt, Ashwin Paranjape, Michele Bevilacqua, Fabio Petroni,
and Percy Liang. Lost in the middle: How language models use long contexts, 2023. URL
https://arxiv.org/abs/2307.03172.
[40] Charlie Snell, Jaehoon Lee, Kelvin Xu, and Aviral Kumar. Scaling llm test-time compute
optimally can be more effective than scaling model parameters, 2024. URL https://arxiv.
org/abs/2408.03314.
[41] Xuezhi Wang, Jason Wei, Dale Schuurmans, Quoc Le, Ed Chi, Sharan Narang, Aakanksha
Chowdhery, and Denny Zhou. Self-consistency improves chain of thought reasoning in lan-
guage models, 2023. URL https://arxiv.org/abs/2203.11171.
15
[42] Alberto Abadie, Susan Athey, Guido W Imbens, and Jeffrey M Wooldridge. When should you
adjust standard errors for clustering?*. The Quarterly Journal of Economics, 138(1):1–35, 10
2022. ISSN 0033-5533. doi: 10.1093/qje/qjac038. URL https://doi.org/10.1093/qje/
qjac038.
[43] A. Colin Cameron and Douglas L. Miller. A practitioner’s guide to cluster-robust inference.
Journal of Human Resources, 50(2):317–372, 2015. ISSN 0022-166X. doi: 10.3368/jhr.50.2.
317. URL https://jhr.uwpress.org/content/50/2/317.
[44] R.M. Bell and Daniel Mccaffrey. Bias reduction in standard errors for linear regression with
multi-stage samples. Survey Methodology, 28:169–181, 01 2002.
[45] Stack Overflow.
2024 developer survey:
Technology - integrated development en-
vironment,
2024.
URL https://survey.stackoverflow.co/2024/technology\#
1-integrated-development-environment. Accessed: June 26, 2025.
16
A
Author contributions
Joel Becker and Nate Rush designed, implemented, and led the project.
Beth Barnes gave feedback and guidance on the project.
David Rein contributed substantially to the writing and framing of the results.
B
Extended Discussion
We do not provide evidence that:
Clarification
AI systems do not currently speed up many
or most software developers
We do not claim that our developers or reposito-
ries represent a majority or plurality of software
development work
AI systems do not speed up individuals or
groups in domains other than software de-
velopment
We only study software development
AI systems in the near future will not speed
up developers in our exact setting
Progress is difficult to predict, and there has been
substantial AI progress over the past five years [2]
There are not ways of using existing AI
systems more effectively to achieve posi-
tive speedup in our exact setting
Cursor does not sample many tokens from
LLMs, it may not use optimal prompting/s-
caffolding,
and
domain/repository-specific
training/finetuning/few-shot learning could yield
positive speedup
Developers with much more experience us-
ing AI systems wouldn’t see speedup
While our developers appear to use AI tools com-
petently, it’s plausible that much more experience
could yield speedup. See Section C.2.8 for dis-
cussion/analysis of the impact of developers’ AI
experience on slowdown.
These developers are not sped up on any
tasks
We estimate that some developers experience
speedup
from
AI—see
Figure
17
for
per-
developer speedup estimates.
Table 2: Potential misconceptions about our work: what our evidence does not demonstrate about
AI and developer productivity.
Potential Misreadings of Results
Given both the importance of understanding AI capabili-
ties/risks, and the diversity of perspectives on these topics, we feel it’s important to forestall potential
misunderstandings or over-generalizations of our results. We list claims that we do not provide evi-
dence for in Table 2.
Literature Comparison
Table 3 compares relevant studies that measure the impact of AI tools on
software developer productivity, along key dimensions that distinguish our results from prior work.
Other relevant literature typically does not investigate the impact of AI systems more capable than
GPT-4 (which as of mid-2025 isn’t itself close to the frontier), and does not (in any individual study)
analyze tasks in the wild, with experienced developers, using outcome measures that correspond
directly with productivity and that are fixed before treatment assignment.
17
Paper
Result
AI ≥GPT-
4?
Non-
synthetic
tasks
Experienced,
high-
familiarity
devs
Fixed
outcome
measure
Peng et al. [10]
↑56\% faster
✗
✗
✗
✓
Weber et al. [18]
↑65\% faster
✗
✗
✗
✓
Cui et al. [17]
↑26\% output
✗
✓
✓
✗
Paradis et al. [11]
↑21\% faster
?
✗
✗
✓
Gambacorta et al. [15]
↑55\% output
✗
✓
✓
✗
Yeverechyahu et al. [16]
↑37\% output
✗
✓
✓
✗
Our study
↓19\% slower
✓
✓
✓
✓
Table 3: Overview of key studies measuring the impact of AI tools on software development produc-
tivity. Paradis et al. [11] does not report the model(s) used internally.
C
Factor Analysis
C.1
Factors driving slowdown
We observe 5 setting-specific factors that contribute to our observed slowdown, summarized in Ta-
ble 1.
C.1.1
Over-optimism about AI usefulness (Direct productivity loss)
On AI-allowed issues, developers are not required to use AI tools—they are instructed to use AI
assistance as much or as little as they would find helpful. Given this, we might expect the percentage
change in issue completion time due to AI to be lower-bounded at 0\%—if a developer is aiming to
be as productive as possible, and they see that they are being slowed down by an AI tool, we’d expect
them to not continue using the AI tool.
However, developers have an overoptimistic picture of how AI affects their productivity, both before
and after they complete issues. Before starting issues, developers forecast that using AI will reduce
completion time by 24\%. After completing issues, developers estimate that using AI reduced their
issue completion time by 20\% on average. This overoptimistic view may lead developers to overuse
AI assistance, despite its negative effect on their productivity.
C.1.2
High developer familiarity with repositories (Raises developer performance)
We might expect that AI assistance is less helpful on tasks where developers are already highly
skilled, for example, because they have completed similar tasks previously, or because they have all
of the knowledge and skills required already.
Before developers complete each issue, they score how much “Prior Task Exposure”, they have
from 1 - 5, where 1 means “Never done this type of task before”, and 5 means “I am an expert at
this type of task”. We consider ratings of less than or equal to 3 to mean they have low prior task
exposure, and ratings higher than 3 to imply high prior task exposure (we selected this cutoff to
achieve roughly balanced ratings).
Figure 7 gives moderate evidence that developers are slowed down more on on issues where they
have high prior task exposure.
Developers also rate (before issues are randomized and completed) “External Resource Needs” from
1 - 3, where 1 means “I need extensive documentation / reference material / googling to complete this
task” and 3 means “I could complete this task entirely from memory and experience.” We consider
ratings of less than or equal to 2 to mean they have low external resource needs, and ratings higher
than 2 to imply high external resource needs (we selected this cutoff to achieve roughly balanced
ratings). Figure 7 presents moderate evidence that developers are slowed down more on issues where
they need fewer external resources.
18
Figure 7: Developers are slowed down more on issues where they self-report having significant
prior task exposure, and on issues where they self-report having low external resource needs (e.g.
documentation, reference materials). We only collected this data for the latter half of issues com-
pleted in the study—this is why we have a smaller number of issues for our baseline slowdown
estimate. See Section D.4 for details on how we estimate heterogeneous treatment effects.
Qualitatively, developers note that AI is particularly helpful when working on unfamiliar issues, and
less helpful when working on familiar ones.
One developer working with unfamiliar datasets found that AI was helpful in answering “general
questions about e.g. EICAR.” Another developer noted that Cursor was “super helpful in figuring
out how to write a [frontend test.] I didn’t know how to do this before and on my third time asking
cursor for help with it, it came up with this solution.” Another developer, working with Git hooks,
noted that “Given that it was my first time with Git hooks, without AI the implementation would’ve
taken me [3 additional hours].” Sometimes, portions of one’s own codebase can be as unknown as
a new API. One developer noted that “cursor found a helper test function that I didn’t even know
existed when I asked it how we tested deprecations.”
On the other hand, developers note that AI is much less helpful on issues where they are expert. One
developer notes that “if I am the dedicated maintainer of a very specialized part of the codebase,
there is no way agent mode can do better than me.”
Broadly, we present moderate evidence that on the issues in our study, developers are slowed down
more when they have high prior task exposure and lower external resource needs. We hypothe-
size that analogously, AI helps our developers less compared to existing literature [10; 11] because
our developers have substantially more experience on their respective repositories (5 years and
1,500 commits on average). This would be consistent with the experience/familiarity effects ob-
served in Noy and Zhang [23]; Cui et al. [17].
C.1.3
Large and complex repositories (Limits AI performance)
Developers qualitatively note LLM tooling performs worse in more complex environments. One
developer says “it also made some weird changes in other parts of the code that cost me time to
find and remove [...] My feeling is the refactoring necessary for this PR was “too big” [and genAI]
introduced as many errors as it fixed.” Another developer comments that one prompt “failed to
properly apply the edits and started editing random other parts of the file,” and that these failures
seemed to be heavily related to “the size of a single file it is attempting to perform edits on.”.
We hypothesize that analogously to these size and complexity effects within our study, AI broadly
helps our developers less compared to existing randomized controlled trials (RCTs) measuring
speedup from AI tools because of the overall size and complexity of the repositories included—
participating repositories are on average about 10 years old and contain >1,100,000 lines of code,
compared to the more greenfield pro
\end{document}
